\documentclass[10.5pt,a4paper]{article}
\usepackage[T1]{fontenc}
\usepackage{charter}
\usepackage[scaled=0.92]{helvet}
\usepackage[a4paper,top=15mm,bottom=15mm,left=17mm,right=17mm]{geometry}
\usepackage{microtype}
\usepackage{enumitem}
\usepackage{titlesec}
\usepackage{tabularx}
\usepackage[dvipsnames]{xcolor}
\usepackage[hidelinks]{hyperref}
\usepackage{fancyhdr}

\definecolor{accent}{HTML}{1F3F7A}
\pagestyle{fancy}\fancyhf{}
\renewcommand{\headrulewidth}{0pt}
\fancyfoot[R]{\footnotesize\color{gray}DongGeon Kwak \quad \thepage}
\setlength{\parindent}{0pt}
\linespread{1.04}

\titleformat{\section}{\color{accent}\sffamily\bfseries\small}{}{0pt}{\MakeUppercase}[\vspace{-6pt}\color{accent}\rule{\textwidth}{0.5pt}]
\titlespacing*{\section}{0pt}{11pt}{6pt}

\setlist[itemize]{leftmargin=13pt,topsep=2pt,itemsep=1.5pt,parsep=0pt,label={\small\textbullet}}

% \entry{Title}{Right}{Subtitle}{Right2}
\newcommand{\entry}[4]{%
  \begin{tabularx}{\textwidth}{@{}X@{\hspace{6pt}}r@{}}
    \textbf{#1} & #2\\[-1pt]
    \textit{#3} & \textit{#4}
  \end{tabularx}\vspace{-2pt}}
\newcommand{\oneline}[2]{%
  \begin{tabularx}{\textwidth}{@{}X@{\hspace{6pt}}r@{}}
    #1 & #2
  \end{tabularx}\vspace{-2pt}}
\newcommand{\gap}{\vspace{5pt}}

\begin{document}

\begin{center}
  {\fontsize{22}{26}\selectfont\bfseries DongGeon (Lloyd) Kwak}\\[5pt]
  {\small
  \href{mailto:kwaklloyd@gmail.com}{kwaklloyd@gmail.com}\enspace$\cdot$\enspace
  \href{https://lloydkwak.github.io}{lloydkwak.github.io}\enspace$\cdot$\enspace
  \href{https://github.com/lloydkwak}{github.com/lloydkwak}\enspace$\cdot$\enspace
  \href{https://www.linkedin.com/in/lloyd-kwak-347a79435}{linkedin.com/in/lloyd-kwak-347a79435}}
\end{center}
\vspace{-2pt}

\section{Research Interests}
Generalist robot learning, cross-embodiment learning, lifelong robot learning (adaptation, failure detection and recovery), and long-horizon manipulation with task and motion planning.

\section{Education}
\entry{Seoul National University of Science and Technology (SeoulTech)}{Seoul, Korea}
      {B.E. in Computer Science and Engineering; Minor in Applied Artificial Intelligence}{Feb.\ 2022 -- Aug.\ 2027 (expected)}
\begin{itemize}
  \item GPA 4.13/4.50 (8th of 114 in major).
  \item A+ in Linear Algebra, Discrete Mathematics, Machine Learning, Computer Vision, Data Analysis, Operating Systems; also completed Probability and Random Variables, Reinforcement Learning.
\end{itemize}

\section{Research Experience}
\entry{Robot Learning Research: FAVOR}{Seoul, Korea}
      {Research feedback: Prof.\ Ji-Hyeong Han; weekly lab meetings}{Jul.\ 2025 -- Present}
\begin{itemize}
  \item \textbf{FAVOR} (first-author manuscript in preparation): training-free correction of a pretrained diffusion policy when a joint of the robot arm locks or loses range.
  \item Compared seven corrections applied at three stages of inference over 112 simulated fault conditions with paired seeds. Retargeting the intended end-effector motion position-first raised mean success from 0.37 to 0.60 (0.11 to 0.48 with a locked joint).
  \item Showed that the end-effector error left after correction predicts which faults are recoverable (AUC 0.94); now validating on a real xArm7.
\end{itemize}
\gap
\entry{MTREA: Multi-Task Retentive Embodiment Adaptation}{SeoulTech}
      {Capstone Design project, two-person team}{Sep.\ 2026 -- Present}
\begin{itemize}
  \item Studying whether fine-tuning the $\pi_{0.5}$ VLA on a low-cost SO-101 arm degrades its performance on held-out tasks; co-designed the protocol (naive fine-tuning baseline, four held-out task families, per-task oracle).
  \item Built calibration and teleoperation for data collection, remote GPU training, and asynchronous policy serving on the real arm.
\end{itemize}

\section{Manuscripts and Patents}
\begin{itemize}[leftmargin=28pt,labelwidth=22pt,labelsep=6pt,align=left,itemsep=3pt]
  \item [M1]\enspace \textbf{D.\ Kwak}. ``FAVOR: Fault-Aware Visuomotor Policy Adaptation via Ordered Retargeting.'' \textit{Manuscript in preparation}, 2026. \href{https://github.com/lloydkwak/FAVOR}{[code]}
  \item [P1]\enspace K.\ Lee, G.\ Lee, \textbf{D.\ Kwak}. ``Real-time crack detection system for construction sites.'' Korean Patent 10-2647135, registered Mar.\ 12, 2024 (filed Oct.\ 31, 2022).
\end{itemize}

\section{Selected Projects}
\entry{MakeYourBrick: One Photo to a Buildable LEGO Model}{May 2026}
      {SeoulTech ST Startup Club Support Program, team PixBricks}{}
\begin{itemize}
  \item Web app that reconstructs an object from a single photo (SAM2, SAM 3D Objects) and converts it into a buildable LEGO model; handled color mapping and shadow artifacts.
\end{itemize}
\gap
\entry{RL-Arena: Competitive Reinforcement-Learning Platform}{Nov.\ -- Dec.\ 2025}
      {Open Source Software course; two-person team; selected as an outstanding project}{}
\begin{itemize}
  \item Platform for submitting RL agents and ranking them by ELO. I built the environment package (published on PyPI), the match executor, and most of the backend.
\end{itemize}
\gap
\entry{Frequency-Aware Preprocessing for Diffusion Image Detection}{Dec.\ 2025}
      {Team project, two people}{}
\begin{itemize}
  \item Found that diffusion-generated images have a narrower pixel-intensity spread and built preprocessing around it; ResNet-50 test accuracy rose from 72.6\% to 96.3\% on 115k images.
\end{itemize}
\gap
\entry{Logisim MLP: A Neural Network Built from Logic Gates}{Spring 2025}
      {Digital Logic and Systems course; selected as an outstanding project}{}
\begin{itemize}
  \item A 256-128-10 MNIST classifier implemented as a digital logic circuit, with trained weights stored in ROM.
\end{itemize}

\section{Work Experience}
\oneline{\textbf{CHASER}, Software Developer}{Nov.\ 2022 -- Jan.\ 2023}
\oneline{\textbf{Laver}, Software Developer}{Jun.\ 2022 -- Oct.\ 2022}
\begin{itemize}
  \item Developed an object-detection system from model to product integration, backend APIs, and an IFC-to-XML converter for building models.
\end{itemize}

\section{Honors}
\oneline{Outstanding term project selections: Digital Logic and Systems, Open Source Software, Computer Vision}{2025 -- 2026}
\oneline{SeoulTech ST Startup Club Support Program, team PixBricks}{2026}

\section{Skills}
\begin{tabularx}{\textwidth}{@{}l@{\hspace{12pt}}X@{}}
  \textbf{Programming}   & Python, C++\\
  \textbf{Robot learning}& PyTorch, LeRobot, MuJoCo, LIBERO\\
  \textbf{Tools}         & Docker, Git, Linux\\
  \textbf{Languages}     & Korean (native); English (TOEIC Speaking, Level AL)\\
  \textbf{Certification} & Craftsman Programming, HRD Korea (2023)
\end{tabularx}

\end{document}
